% ============================================================================
% CHAPTER 3: WORD REPRESENTATIONS AND EMBEDDINGS
% The foundation of all modern NLP --- from one-hot vectors to contextual
% embeddings.  Understanding this evolution is tested in every NLP interview.
% ============================================================================

\chapter{Word Representations and Embeddings}
\label{chap:word_representations}

\begin{tldr}
Word representations are the foundation of modern NLP---the quality of your
embeddings determines the ceiling of your system's performance.  This chapter
traces the evolution from one-hot vectors through Word2Vec, GloVe, and FastText
to contextual embeddings (ELMo) and modern sentence embeddings (SBERT, E5).
Understanding this evolution is critical for interviews because it explains
\emph{why} each subsequent approach was developed, what problem it solved, and
what limitations remained.  Every FAANG NLP interview will test at least one
concept from this chapter, and candidates who can derive the Word2Vec objective
and explain the distributional hypothesis from first principles earn
significantly stronger signals than those who jump straight to ``just use
BERT.''
\end{tldr}

\bigskip

% ============================================================================
% SECTION 1: THE REPRESENTATION PROBLEM
% ============================================================================
\section{The Representation Problem}
\label{sec:representation_problem}

Before any machine learning model can process text, words must be converted
into numerical representations.  The quality and structure of these
representations fundamentally constrain what the model can learn.  This section
examines why the naive approach fails and how the field arrived at
distributional representations.

\subsection{One-Hot Encoding: The Naive Baseline}

The simplest representation assigns each word in the vocabulary a unique index
and encodes it as a vector with a single 1 in that position and 0s everywhere
else.  For a vocabulary of size $V$:
\[
    \text{``cat''} = [0, 0, \ldots, 1, \ldots, 0]^\top \in \R^V
\]

One-hot encoding has three fatal flaws:

\begin{enumerate}[label=\textbf{\arabic*.}]
    \item \textbf{No notion of similarity.}  For any two distinct words $w_i$
    and $w_j$, the inner product of their one-hot vectors is zero:
    $\inner{\ve_{w_i}}{\ve_{w_j}} = 0$.  ``Cat'' is equally distant from
    ``kitten'' as it is from ``refrigerator.''

    \item \textbf{Dimensionality explosion.}  Production vocabularies range from
    30K to 250K tokens.  Representing each word as a 100K-dimensional sparse
    vector is wasteful and computationally expensive for downstream models.

    \item \textbf{No generalization.}  A model that learns something about
    ``cat'' gains no information about ``kitten,'' because the representations
    share no structure.  Every word is an island.
\end{enumerate}

The fundamental insight that unlocked progress was that we do not need to
represent words as atomic symbols.  Instead, we can represent them as
\emph{dense, low-dimensional vectors} in which geometric proximity encodes
semantic similarity.

\subsection{The Distributional Hypothesis}

\begin{keyinsight}[The Foundation of All Embeddings]
\emph{``A word is characterized by the company it keeps.''}
--- J.~R.~Firth, 1957

The distributional hypothesis states that words appearing in similar contexts
tend to have similar meanings.  This single idea is the theoretical foundation
for \emph{every} word embedding method---from count-based approaches (PMI, SVD)
to predictive models (Word2Vec, GloVe) to contextual embeddings (ELMo, BERT).
If an interviewer asks ``why do word embeddings work?'' this is the answer.
\end{keyinsight}

The hypothesis is operationalized differently across methods, but the core
principle is always the same: extract statistical patterns from word
co-occurrences in large corpora, and use those patterns to define the
representation.

\subsection{Count-Based Methods: Co-occurrence and PMI}

Before neural embeddings, researchers built word representations from
co-occurrence statistics.

\paragraph{Co-occurrence Matrix.}
Given a corpus, define a window size $L$ and count how often each pair of words
$(w_i, w_j)$ appears within $L$ tokens of each other.  This produces a
$V \times V$ matrix $\mX$ where $X_{ij}$ is the co-occurrence count.

\paragraph{Pointwise Mutual Information (PMI).}
Raw counts are dominated by frequent words.  PMI re-weights co-occurrences:
\[
    \text{PMI}(w_i, w_j)
    = \log \frac{P(w_i, w_j)}{P(w_i)\, P(w_j)}
    = \log \frac{X_{ij} \cdot N}{X_{i*}\, X_{*j}}
\]
where $N = \sum_{i,j} X_{ij}$ is the total count, $X_{i*} = \sum_j X_{ij}$,
and $X_{*j} = \sum_i X_{ij}$.  In practice, negative values are clipped to
zero (\textbf{Positive PMI} or PPMI), since negative PMI is unreliable for
rare events.

\paragraph{Dimensionality Reduction via SVD.}
The PPMI matrix is $V \times V$ and sparse.  Applying truncated Singular Value
Decomposition (SVD) to retain the top $d$ singular values produces dense
$d$-dimensional word vectors:
\[
    \mX_{\text{PPMI}} \approx \mU_d \bm{\Sigma}_d \mV_d^\top
\]
The rows of $\mU_d \bm{\Sigma}_d^{1/2}$ (or simply $\mU_d$) serve as word
vectors.  Levy and Goldberg (2014) showed that this count-based approach
produces embeddings competitive with Word2Vec, establishing the deep connection
between count-based and predictive methods.

\begin{warningbox}
A common interview mistake is treating count-based methods as ``ancient
history'' unrelated to modern embeddings.  In fact, GloVe explicitly bridges
count-based and predictive approaches, and understanding this connection
demonstrates genuine depth.  Interviewers at Google and Meta have been known
to ask: ``Can you explain how SVD on a PPMI matrix relates to Word2Vec?''
\end{warningbox}


% ============================================================================
% SECTION 2: WORD2VEC
% ============================================================================
\section{Word2Vec --- Know This Cold}
\label{sec:word2vec}

Word2Vec (Mikolov et al., 2013) transformed NLP by showing that simple neural
networks trained on large text corpora produce word vectors with remarkable
properties, including the now-famous linear analogy structure.  It remains one
of the most frequently tested topics in NLP interviews because it tests
mathematical reasoning, optimization understanding, and awareness of practical
training tricks.

\subsection{Skip-Gram: Predict Context from Center Word}

The Skip-gram model defines the following task: given a center word $w_I$,
predict each context word $w_O$ within a window of size $L$.

\paragraph{Architecture.}
Each word $w$ has two learned vectors: an \emph{input embedding}
$\vw_w \in \R^d$ (used when $w$ is the center word) and an \emph{output
embedding} $\vw'_w \in \R^d$ (used when $w$ is a context word).  The model
has no hidden layers---it is a shallow architecture with a single projection.

\begin{figure}[H]
\centering
\begin{tikzpicture}[
    node distance=1.2cm and 2.5cm,
    box/.style={rectangle, draw=deepblue, fill=lightblue, minimum width=2.0cm,
                minimum height=0.7cm, font=\small\sffamily, rounded corners=2pt,
                thick},
    vec/.style={rectangle, draw=deepblue!60, fill=white, minimum width=1.4cm,
                minimum height=0.7cm, font=\small\sffamily, rounded corners=2pt},
    arr/.style={-{Stealth[length=5pt]}, thick, deepblue!80},
    lbl/.style={font=\scriptsize\sffamily, text=gray!70!black}
]

% Input
\node[box] (input) {$w_I$};
\node[lbl, below=0.1cm of input] {Center word};

% Projection
\node[box, right=2.5cm of input, fill=deepblue!12] (proj) {$\vw_{w_I} \in \R^d$};
\node[lbl, below=0.1cm of proj] {Projection (lookup)};

% Output words
\node[vec, right=2.5cm of proj, yshift=1.8cm]  (o1) {$w_{-L}$};
\node[vec, right=2.5cm of proj, yshift=0.6cm]  (o2) {$w_{-1}$};
\node[vec, right=2.5cm of proj, yshift=-0.6cm] (o3) {$w_{+1}$};
\node[vec, right=2.5cm of proj, yshift=-1.8cm] (o4) {$w_{+L}$};

% Dots
\node[font=\large, text=gray] at ($(o2)!0.5!(o3)$) {$\vdots$};

\node[lbl, below=0.1cm of o4] {Context words};

% Arrows
\draw[arr] (input) -- (proj) node[midway, above, lbl] {$\mW \in \R^{V \times d}$};
\draw[arr] (proj) -- (o1) node[midway, above, lbl, sloped] {$\softmax$};
\draw[arr] (proj) -- (o2);
\draw[arr] (proj) -- (o3);
\draw[arr] (proj) -- (o4) node[midway, below, lbl, sloped] {$\softmax$};

\end{tikzpicture}
\caption{Skip-gram architecture.  The center word is projected to a
$d$-dimensional embedding, and the model predicts each context word within a
window of size $L$ independently.}
\label{fig:skipgram}
\end{figure}

\paragraph{Probability Model.}
The probability of observing context word $w_O$ given center word $w_I$ is
defined via the softmax over the entire vocabulary:

\begin{mathresult}{Skip-Gram Softmax Objective}
\[
    P(w_O \given w_I)
    = \frac{\exp\!\bigl(\vw'_{w_O}{}^\top \vw_{w_I}\bigr)}
           {\displaystyle\sum_{w=1}^{V}
            \exp\!\bigl(\vw'_{w}{}^\top \vw_{w_I}\bigr)}
\]
where $\vw_{w_I} \in \R^d$ is the input embedding of center word $w_I$,
$\vw'_{w_O} \in \R^d$ is the output embedding of context word $w_O$, and
$V$ is the vocabulary size.
\end{mathresult}

\paragraph{Training Objective.}
Given a corpus of words $w_1, w_2, \ldots, w_T$, maximize the average
log-probability of context words:
\[
    \loss_{\text{skip-gram}}
    = \frac{1}{T} \sum_{t=1}^{T}
      \sum_{\substack{-L \leq j \leq L \\ j \neq 0}}
      \log P(w_{t+j} \given w_t)
\]

The summation over $j$ iterates over all context positions within the window.
The full softmax requires computing the partition function (denominator) over
all $V$ words, which is $O(V)$ per training example---prohibitively expensive
for vocabularies of 100K+ words.

\subsection{CBOW: Predict Center from Context}

The Continuous Bag-of-Words (CBOW) model reverses the prediction direction:
given all context words within the window, predict the center word.  The
context is represented as the average of context word embeddings:
\[
    \bar{\vw} = \frac{1}{2L} \sum_{\substack{-L \leq j \leq L \\ j \neq 0}}
                \vw_{w_{t+j}}
\]
\[
    P(w_t \given \text{context})
    = \frac{\exp(\vw'_{w_t}{}^\top \bar{\vw})}
           {\sum_{w=1}^{V} \exp(\vw'_w{}^\top \bar{\vw})}
\]

\paragraph{CBOW vs.\ Skip-gram in Practice.}
\begin{itemize}
    \item \textbf{CBOW} is faster to train (one prediction per window, not
    $2L$) and produces better embeddings for \emph{frequent} words (averages
    context, smoothing out noise).
    \item \textbf{Skip-gram} produces better embeddings for \emph{rare} words
    and smaller datasets (each context word provides a separate training
    signal, so rare words get more updates).
\end{itemize}

In practice, Skip-gram with negative sampling became the dominant choice and
is the version interviewers expect you to know in detail.

\subsection{Negative Sampling: Making Training Tractable}

The full softmax is $O(V)$ per example.  Negative sampling reformulates the
problem from a $V$-class classification to a set of binary classifications.

\paragraph{Core Idea.}
Instead of predicting the correct context word from the entire vocabulary,
train a binary classifier to distinguish \emph{real} (center, context) pairs
from \emph{noise} pairs where the context word is drawn from a noise
distribution.

\begin{mathresult}{Negative Sampling (NEG) Objective}
\[
    \loss_{\text{NEG}}
    = -\log \sigma\!\bigl(\vw'_{w_O}{}^\top \vw_{w_I}\bigr)
      - \sum_{i=1}^{k} \E_{w_i \sim P_n(w)} \Bigl[
        \log \sigma\!\bigl(-\vw'_{w_i}{}^\top \vw_{w_I}\bigr)
      \Bigr]
\]
where $\sigma(x) = 1/(1+e^{-x})$ is the sigmoid function, $k$ is the number
of negative samples, and $P_n(w)$ is the noise distribution.
\end{mathresult}

\paragraph{Why This Works.}
The first term pushes the dot product of the true (center, context) pair
toward $+\infty$ (high similarity).  The second term pushes the dot products
of the noise pairs toward $-\infty$ (low similarity).  NEG is a
simplification of noise-contrastive estimation (NCE): unlike NCE, it does
\emph{not} recover the full softmax gradient, and it should not be confused
with the mutual-information bound maximized by InfoNCE (Oord et al., 2018).
What NEG actually optimizes is characterized by Levy and Goldberg (2014):
at its optimum, it implicitly factorizes a shifted pointwise mutual
information matrix (see the shifted-PMI result later in this section).

\paragraph{Noise Distribution.}
The choice of $P_n(w)$ matters significantly.  Mikolov et al.\ found that the
\emph{unigram distribution raised to the 3/4 power} works best:
\[
    P_n(w) = \frac{f(w)^{3/4}}{\sum_{w'} f(w')^{3/4}}
\]
where $f(w)$ is the frequency of word $w$.  The 3/4 exponent smooths the
distribution, increasing the probability of sampling moderately rare words
relative to the raw unigram distribution.  This makes the negative samples
more challenging and informative.

\paragraph{Number of Negative Samples.}
Mikolov et al.\ recommend $k = 5$--$20$ for small datasets and $k = 2$--$5$
for large ones.  More negatives provide a better approximation of the full
softmax but increase computation.

\subsection{Subsampling Frequent Words}

Extremely frequent words like ``the,'' ``a,'' and ``is'' provide little
information---co-occurring with ``the'' tells you almost nothing about a
word's meaning.  Subsampling randomly discards occurrences of frequent words
during training with probability:
\[
    P(\text{discard } w) = 1 - \sqrt{\frac{t}{f(w)}}
\]
where $f(w)$ is the word's frequency and $t$ is a threshold (typically
$10^{-5}$).  Words with frequency above $t$ are increasingly likely to be
discarded.  This both speeds up training and improves embedding quality for
rare words.

\subsection{Properties of Word2Vec Embeddings}

The most celebrated property is \textbf{linear analogy structure}:
\[
    \vw_{\text{king}} - \vw_{\text{man}} + \vw_{\text{woman}}
    \approx \vw_{\text{queen}}
\]

More precisely, the analogy ``man is to king as woman is to ?'' is solved by:
\[
    \argmax_{w} \cos(\vw_w, \;\vw_{\text{king}} - \vw_{\text{man}}
                               + \vw_{\text{woman}})
\]

This linear structure arises because the training objective implicitly
factorizes a shifted PMI matrix (Levy and Goldberg, 2014):
\[
    \vw_{w_I}{}^\top \vw'_{w_O}
    \approx \text{PMI}(w_I, w_O) - \log k
\]
where $k$ is the number of negative samples.

Other important properties:
\begin{itemize}
    \item \textbf{Semantic clusters:} Similar words cluster together
    (countries, animals, professions form coherent regions).
    \item \textbf{Syntactic regularities:} ``walking'' $-$ ``walked'' $+$
    ``swimming'' $\approx$ ``swam'' (tense analogies).
    \item \textbf{Compositionality (limited):} Simple vector addition
    sometimes captures phrase meaning, but this breaks down for complex
    constructions.
\end{itemize}

\subsection{Hyperparameter Guide}

\begin{table}[H]
\centering
\caption{Word2Vec hyperparameters and their effects.}
\label{tab:w2v_hyperparams}
\begin{tabularx}{\textwidth}{l l X}
\toprule
\textbf{Hyperparameter} & \textbf{Typical Range} & \textbf{Effect} \\
\midrule
Dimension $d$ & 100--300 & Higher $d$ captures more nuance but
    requires more data to train well. 300 is a common default. \\
Window size $L$ & 2--10 & Smaller windows ($\leq 3$) capture syntactic
    relationships (POS, morphology). Larger windows ($\geq 5$) capture
    topical/semantic relationships. \\
Negative samples $k$ & 2--20 & More negatives = better softmax
    approximation but slower training. 5 is a standard choice. \\
Min count & 5--100 & Discard words appearing fewer than this many times.
    Reduces vocabulary and noise from rare words. \\
Subsampling $t$ & $10^{-5}$--$10^{-3}$ & Lower = more aggressive
    downsampling of frequent words. \\
Learning rate & 0.01--0.025 & Linearly decayed during training. \\
\bottomrule
\end{tabularx}
\end{table}

\subsection{Limitations of Word2Vec}

\begin{enumerate}[label=\textbf{\arabic*.}]
    \item \textbf{Static representations.}  Each word has exactly one vector
    regardless of context.  ``Bank'' in ``river bank'' and ``bank account''
    maps to the same point---a fundamental limitation that contextual
    embeddings (Section~\ref{sec:contextual_embeddings}) address.

    \item \textbf{Out-of-vocabulary (OOV) words.}  Any word not seen during
    training has no representation.  This is a serious problem for
    morphologically rich languages and specialized domains.  FastText
    (Section~\ref{sec:fasttext}) solves this.

    \item \textbf{No subword information.}  Word2Vec treats each word as an
    atomic unit.  It cannot exploit the fact that ``unhappiness,''
    ``unhappy,'' and ``happiness'' share morphemes.

    \item \textbf{Frequency bias.}  Despite subsampling, rare words have
    lower-quality embeddings because they receive fewer training updates.

    \item \textbf{No sentence-level representation.}  Word2Vec produces
    word vectors, not sentence or document vectors (though averaging provides
    a surprisingly decent baseline).
\end{enumerate}

\begin{warningbox}
Candidates often describe Word2Vec as ``just a lookup table'' or dismiss it as
``replaced by BERT.''  This is a red flag.  Interviewers expect you to know
the training objective, explain negative sampling, derive the connection to
PMI, and discuss the hyperparameter tradeoffs.  Word2Vec's ideas (projection
layers, contrastive training, subsampling) are directly ancestral to modern
approaches.
\end{warningbox}


% ============================================================================
% SECTION 3: GLOVE
% ============================================================================
\section{GloVe --- Global Vectors for Word Representation}
\label{sec:glove}

GloVe (Pennington et al., 2014) was designed to bridge the gap between
count-based methods (SVD on co-occurrence matrices) and predictive methods
(Word2Vec).  Its key insight is that the \emph{ratios} of co-occurrence
probabilities encode meaning more directly than raw probabilities.

\subsection{Motivation: Co-occurrence Ratios}

Consider three words: ``ice,'' ``steam,'' and ``solid.''  Let
$P(k \given w)$ denote the probability that word $k$ appears in the context
of word $w$.  The ratio $P(k \given \text{ice}) / P(k \given \text{steam})$
is revealing:

\begin{table}[H]
\centering
\begin{tabular}{lcccc}
\toprule
& \textbf{$k$ = solid} & \textbf{$k$ = gas} & \textbf{$k$ = water}
& \textbf{$k$ = fashion} \\
\midrule
$P(k \given \text{ice})$ & Large & Small & Large & Small \\
$P(k \given \text{steam})$ & Small & Large & Large & Small \\
\textbf{Ratio} & \textbf{Large} & \textbf{Small} & $\approx 1$ & $\approx 1$ \\
\bottomrule
\end{tabular}
\end{table}

Words that discriminate between ``ice'' and ``steam'' (like ``solid'' and
``gas'') produce extreme ratios, while words related to both (``water'') or
neither (``fashion'') produce ratios near 1.  GloVe's objective is designed
so that the learned vectors capture these ratio patterns.

\subsection{The GloVe Objective}

GloVe seeks word vectors such that their dot product approximates the
logarithm of the co-occurrence count:

\begin{mathresult}{GloVe Objective}
\[
    \loss_{\text{GloVe}}
    = \sum_{i,j=1}^{V} f(X_{ij})\,
      \bigl(\vw_i^\top \vw'_j + b_i + b'_j - \log X_{ij}\bigr)^2
\]
where $X_{ij}$ is the co-occurrence count between words $i$ and $j$,
$\vw_i, \vw'_j \in \R^d$ are word vectors, $b_i, b'_j$ are scalar biases,
and $f(\cdot)$ is a weighting function.
\end{mathresult}

\paragraph{The Weighting Function $f(x)$.}
The weighting function prevents extremely frequent co-occurrences from
dominating the loss:
\[
    f(x) = \begin{cases}
        (x / x_{\max})^\alpha & \text{if } x < x_{\max} \\
        1 & \text{otherwise}
    \end{cases}
\]
Pennington et al.\ use $x_{\max} = 100$ and $\alpha = 3/4$.  This ensures
that (a) zero co-occurrences are not penalized (since $f(0) = 0$), and (b)
very high co-occurrences are capped at weight 1.

\subsection{Why GloVe Works}

The objective can be derived from the requirement that the dot product of word
vectors encode co-occurrence probability ratios.  Starting from the
desideratum:
\[
    F\bigl((\vw_i - \vw_j)^\top \vw'_k\bigr)
    = \frac{P(k \given i)}{P(k \given j)}
\]
and requiring $F$ to be an exponential (to map addition in vector space to
multiplication in probability space), one arrives at:
\[
    \vw_i^\top \vw'_k = \log P(k \given i) = \log X_{ik} - \log X_i
\]
The bias terms absorb the $\log X_i$ normalization, yielding the objective
above.

\subsection{GloVe vs.\ Word2Vec}

\begin{table}[H]
\centering
\caption{Comparison of Word2Vec and GloVe.}
\label{tab:w2v_vs_glove}
\begin{tabularx}{\textwidth}{l X X}
\toprule
\textbf{Aspect} & \textbf{Word2Vec (Skip-gram + NEG)} & \textbf{GloVe} \\
\midrule
Training data & Local context windows (online, streaming) &
    Global co-occurrence matrix (requires pre-computation) \\
Objective & Predictive (maximize log-probability of context words) &
    Reconstructive (minimize weighted squared error on log counts) \\
Global statistics & Implicitly captured through many passes over data &
    Explicitly captured in the co-occurrence matrix \\
Scalability & Streams data; handles very large corpora &
    Requires $O(V^2)$ co-occurrence matrix (sparse, but large) \\
Performance & Comparable on analogy tasks &
    Comparable on analogy tasks; sometimes better on similarity \\
Connection & Implicitly factorizes a shifted PMI matrix &
    Explicitly optimizes a log co-occurrence factorization \\
\bottomrule
\end{tabularx}
\end{table}

\begin{keyinsight}[The Unification of Count-Based and Predictive Methods]
Levy and Goldberg (2014) showed that Word2Vec's Skip-gram with negative
sampling implicitly factorizes a shifted PMI matrix.  GloVe directly optimizes
a weighted factorization of the log co-occurrence matrix.  Both are
fundamentally performing matrix factorization on co-occurrence statistics---the
difference is in \emph{how} they weight different entries and \emph{how} they
optimize.  This insight, frequently tested in research-track interviews,
shows that the predictive vs.\ count-based distinction is less fundamental
than it appears.
\end{keyinsight}


% ============================================================================
% SECTION 4: FASTTEXT
% ============================================================================
\section{FastText --- Subword Embeddings}
\label{sec:fasttext}

FastText (Bojanowski et al., 2017) extends Word2Vec by representing each word
as a \emph{bag of character n-grams}, enabling the model to capture
morphological structure and handle out-of-vocabulary words.

\subsection{Subword Representation}

Each word $w$ is decomposed into a set of character n-grams $\mathcal{G}_w$
(typically $n = 3$ to $6$), plus special boundary markers \texttt{<} and
\texttt{>} and the word itself:

\paragraph{Example.}
For the word ``where'' with $n \in \{3, 4, 5\}$:
\[
    \mathcal{G}_{\text{where}} = \{\texttt{<wh},\; \texttt{whe},\;
    \texttt{her},\; \texttt{ere},\; \texttt{re>},\; \texttt{<whe},\;
    \texttt{wher},\; \texttt{here},\; \texttt{ere>},\; \texttt{<wher},\;
    \texttt{where},\; \texttt{here>},\; \texttt{<where>}\}
\]

The word's embedding is the sum of its n-gram embeddings plus a dedicated
whole-word vector (if the word appeared in training):
\[
    \vw_w = \vw_w^{(\text{word})} + \sum_{g \in \mathcal{G}_w} \vw_g
\]

\subsection{Handling Out-of-Vocabulary Words}

FastText's key advantage over Word2Vec: for any unseen word, decompose it into
character n-grams and sum the corresponding vectors.  As long as the n-grams
were observed during training (which is almost always the case), the model
produces a meaningful representation.

This is especially valuable for:
\begin{itemize}
    \item \textbf{Morphologically rich languages} (Turkish, Finnish, German):
    compound words and agglutination create enormous vocabularies.
    ``Donaudampfschifffahrtsgesellschaft'' (German: Danube steamship company)
    is handled naturally by its constituent n-grams.
    \item \textbf{Misspellings and typos:} ``teh'' shares most n-grams with
    ``the,'' so its vector is close.
    \item \textbf{Domain-specific terms:} Medical or legal terminology often
    follows morphological patterns (``-itis,'' ``-ectomy'') that n-grams
    capture.
\end{itemize}

\subsection{FastText in Production}

FastText remains widely used in production systems for several reasons:
\begin{enumerate}
    \item \textbf{Speed:} Training and inference are extremely fast---embedding
    lookup is a simple sum of pre-computed vectors.
    \item \textbf{Lightweight:} No GPU required for inference.
    \item \textbf{Supervised mode:} FastText includes a supervised
    classification mode that is competitive with deep learning for simple text
    classification tasks at a fraction of the compute cost.
    \item \textbf{Multilingual:} Pre-trained vectors are available for 157
    languages.
\end{enumerate}

\begin{table}[H]
\centering
\caption{Comparison of static embedding methods.}
\label{tab:static_embeddings}
\begin{tabularx}{\textwidth}{l c c c X}
\toprule
\textbf{Method} & \textbf{Year} & \textbf{OOV} & \textbf{Morphology}
& \textbf{Key Innovation} \\
\midrule
Word2Vec & 2013 & No & No & Scalable predictive embeddings; negative sampling \\
GloVe & 2014 & No & No & Weighted factorization of global co-occurrence \\
FastText & 2017 & Yes & Yes & Character n-gram decomposition \\
\bottomrule
\end{tabularx}
\end{table}


% ============================================================================
% SECTION 5: CONTEXTUAL EMBEDDINGS
% ============================================================================
\section{Contextual Embeddings --- The Bridge to Modern NLP}
\label{sec:contextual_embeddings}

Static embeddings assign one fixed vector per word.  Contextual embeddings
assign a \emph{different} vector depending on the surrounding context.  This
section covers ELMo, the pivotal model that demonstrated the power of
pre-trained contextual representations and directly paved the way for BERT.

\subsection{Why Context Matters: The Polysemy Problem}

Consider the word ``bank'' in two sentences:
\begin{enumerate}
    \item ``I deposited money at the \textbf{bank}.''
    \item ``We sat on the river \textbf{bank}.''
\end{enumerate}

Word2Vec, GloVe, and FastText all produce the \emph{same} vector for ``bank''
in both sentences---an average of its financial and geographical senses,
which accurately represents neither.  A contextual embedding model produces
\emph{different} vectors: one near ``finance'' and ``deposit,'' the other near
``river'' and ``shore.''

\subsection{ELMo: Embeddings from Language Models}

ELMo (Peters et al., 2018) generates contextual embeddings using a
\emph{bidirectional language model} (biLM).

\paragraph{Architecture.}
ELMo trains a two-layer bidirectional LSTM language model:
\begin{itemize}
    \item \textbf{Forward LM:} Predicts the next token given the left context:
    $P(w_t \given w_1, \ldots, w_{t-1})$.
    \item \textbf{Backward LM:} Predicts the previous token given the right
    context: $P(w_t \given w_{t+1}, \ldots, w_T)$.
    \item The forward and backward LMs share the token embedding layer
    (character CNN) but have separate LSTM parameters.
\end{itemize}

\paragraph{Generating Contextual Representations.}
For a given token $w_t$ in context, ELMo computes $2L + 1$ representations
(where $L$ is the number of LSTM layers):
\begin{itemize}
    \item Layer 0: the context-independent character-based embedding
    (from a character CNN)
    \item Layers 1 through $L$: the concatenation of forward and backward LSTM
    hidden states at each layer:
    $\vh_t^{(\ell)} = [\overrightarrow{\vh}_t^{(\ell)} \,;\,
                       \overleftarrow{\vh}_t^{(\ell)}]$
\end{itemize}

The final ELMo representation is a \emph{task-specific weighted sum} across
layers:
\[
    \text{ELMo}_t = \gamma \sum_{\ell=0}^{L} s_\ell \, \vh_t^{(\ell)}
\]
where $s_\ell$ are learned softmax-normalized weights and $\gamma$ is a
scalar scaling factor.  Both are learned by the \emph{downstream task model}
while the biLM itself stays \emph{frozen}---ELMo is feature-based transfer,
not fine-tuning.

\subsection{What Different Layers Capture}

A key finding from Peters et al.\ is that different layers of the biLM
capture different types of linguistic information:

\begin{table}[H]
\centering
\caption{Information captured at different ELMo layers.}
\label{tab:elmo_layers}
\begin{tabularx}{\textwidth}{l X X}
\toprule
\textbf{Layer} & \textbf{Information Type} & \textbf{Evidence} \\
\midrule
Layer 0 (char CNN) & Morphological, orthographic &
    Good for POS tagging; captures word shape \\
Layer 1 (lower LSTM) & Syntactic &
    Best for tasks like POS tagging, constituency parsing \\
Layer 2 (upper LSTM) & Semantic, word sense disambiguation &
    Best for tasks like NER, sentiment, textual entailment \\
\bottomrule
\end{tabularx}
\end{table}

This finding---that lower layers capture syntax and upper layers capture
semantics---has been consistently replicated in Transformer models (BERT,
GPT) and is a frequently asked interview question.

\begin{keyinsight}[ELMo: The Proof of Concept for Pre-training]
ELMo demonstrated three things that shaped all subsequent NLP research:
(1)~pre-training on unlabeled text produces representations useful for many
downstream tasks, (2)~different layers of a deep model capture different
linguistic phenomena, and (3)~contextual representations dramatically
outperform static embeddings.  BERT built on all three insights, replacing
the LSTM with a Transformer and the feature-based approach with fine-tuning.
\end{keyinsight}


% ============================================================================
% SECTION 6: SENTENCE AND DOCUMENT EMBEDDINGS
% ============================================================================
\section{Sentence and Document Embeddings}
\label{sec:sentence_embeddings}

While word embeddings represent individual tokens, many NLP applications
require fixed-size representations of sentences, paragraphs, or entire
documents.  This section covers the methods from simple baselines to modern
production embedding models.

\subsection{Simple Baselines (Surprisingly Strong)}

\paragraph{Average Word Embeddings.}
The simplest approach: average the word vectors of all tokens in a sentence.
\[
    \vw_{\text{sent}} = \frac{1}{n} \sum_{i=1}^{n} \vw_{w_i}
\]
Despite its simplicity, this often outperforms more complex unsupervised
methods.  Arora et al.\ (2017) showed that computing a \emph{weighted}
average (by smooth inverse frequency, SIF) followed by removing the first
principal component yields competitive results on semantic similarity tasks.

\paragraph{TF-IDF Weighted Average.}
Weight each word's contribution by its TF-IDF score, giving more influence to
discriminative terms:
\[
    \vw_{\text{sent}} = \frac{\sum_{i=1}^{n}
                               \text{tfidf}(w_i) \cdot \vw_{w_i}}
                              {\sum_{i=1}^{n} \text{tfidf}(w_i)}
\]

\subsection{Sentence-BERT (SBERT)}

Reimers and Gurevych (2019) introduced Sentence-BERT, which fine-tunes BERT
using a \emph{siamese/triplet network} to produce semantically meaningful
sentence embeddings.

\paragraph{Architecture.}
Two identical BERT encoders (shared weights) process two sentences
independently.  Mean pooling over token outputs produces fixed-size
sentence embeddings.  The objective depends on the training data:

\begin{itemize}
    \item \textbf{NLI training (classification):}  For sentence pairs with
    labels (entailment, contradiction, neutral), concatenate embeddings
    $[\mathbf{u}; \mathbf{v}; |\mathbf{u} - \mathbf{v}|]$ and pass through a
    softmax classifier.
    \item \textbf{STS training (regression):}  For sentence pairs with
    similarity scores, minimize MSE between cosine similarity and the target
    score.
\end{itemize}

\paragraph{Why SBERT Matters.}
BERT alone is impractical for large-scale semantic search: computing the
similarity of 10,000 sentences requires $\binom{10000}{2} \approx 50$M
cross-encoder forward passes.  SBERT produces pre-computed embeddings,
reducing this to 10,000 forward passes plus simple cosine similarity lookups.

\subsection{SimCSE: Contrastive Learning for Sentence Embeddings}

Gao et al.\ (2021) proposed SimCSE, which uses contrastive learning to produce
high-quality sentence embeddings with a remarkably simple technique.

\paragraph{Unsupervised SimCSE.}
Pass the \emph{same} sentence through the encoder \emph{twice} with different
dropout masks.  The two outputs form a positive pair.  Other sentences in the
mini-batch form negatives.  The loss is InfoNCE:

\begin{mathresult}{SimCSE Unsupervised Contrastive Loss}
\[
    \loss = -\log \frac{e^{\text{sim}(\vh_i, \vh_i^+)/\tau}}
                       {\sum_{j=1}^{N} e^{\text{sim}(\vh_i, \vh_j^+)/\tau}}
\]
where $\vh_i$ and $\vh_i^+$ are embeddings of the same sentence with
different dropout masks, $\text{sim}$ is cosine similarity, $\tau$ is the
temperature, and $N$ is the batch size.
\end{mathresult}

\paragraph{Why Dropout Works as Augmentation.}
Dropout randomly zeroes a fraction of hidden units at each forward pass.  Two
passes of the same input produce \emph{slightly different} representations
that are natural ``positive pairs''---semantically identical but numerically
distinct.  This is a minimal data augmentation that avoids the pitfalls of
text-specific augmentations (e.g., synonym replacement can change meaning).

\paragraph{Supervised SimCSE.}
Uses NLI data: entailment pairs are positives, contradiction pairs are hard
negatives.  This consistently outperforms the unsupervised variant.

\subsection{Modern Embedding Models}

The current generation of sentence embedding models builds on the insights
from SBERT and SimCSE, scaled with more data, better training recipes, and
instruction tuning.

\begin{table}[H]
\centering
\caption{Modern sentence embedding models.}
\label{tab:modern_embeddings}
\begin{tabularx}{\textwidth}{l l l X}
\toprule
\textbf{Model} & \textbf{Org} & \textbf{Year} & \textbf{Key Innovation} \\
\midrule
E5 & Microsoft & 2022 & ``query:'' and ``passage:'' prefixes; trained
    with contrastive loss on diverse paired data \\
BGE & BAAI & 2023 & Instruction-tuned embeddings; strong MTEB
    performance \\
GTE & Alibaba & 2023 & Multi-stage contrastive training; competitive
    at small model sizes \\
Nomic Embed & Nomic AI & 2024 & Long-context (8192 tokens); open-source;
    Matryoshka training \\
\bottomrule
\end{tabularx}
\end{table}

\subsection{Matryoshka Embeddings}

Kusupati et al.\ (2022) introduced Matryoshka Representation Learning (MRL),
where a single model is trained to produce useful embeddings at \emph{multiple}
dimensions simultaneously.  The training loss sums over multiple truncation
points:
\[
    \loss_{\text{MRL}} = \sum_{d \in \mathcal{D}}
                         \loss_{\text{task}}(\vw_{[:d]})
\]
where $\mathcal{D} = \{32, 64, 128, 256, 512, 768\}$ (for example) and
$\vw_{[:d]}$ is the first $d$ dimensions of the embedding.

This is valuable in production: you can use 64-dimensional embeddings for
coarse filtering (fast ANN search) and full 768-dimensional embeddings for
fine-grained re-ranking, all from a single model.

\subsection{Training Objectives for Embedding Models}

\paragraph{Contrastive Loss (InfoNCE / NT-Xent).}
The standard loss for embedding models:
\[
    \loss = -\frac{1}{N} \sum_{i=1}^{N} \log
    \frac{e^{\text{sim}(\vh_i, \vh_i^+)/\tau}}
         {\sum_{j=1}^{N}\bigl(
           e^{\text{sim}(\vh_i, \vh_j^+)/\tau}
           + e^{\text{sim}(\vh_i, \vh_j^-)/\tau}\bigr)}
\]

\paragraph{Hard Negative Mining.}
The quality of negatives is critical for embedding models.  Strategies include:
\begin{itemize}
    \item \textbf{In-batch negatives:} Other examples in the mini-batch serve
    as negatives.  Simple but effective; requires large batch sizes.
    \item \textbf{BM25 negatives:} Documents that score highly on BM25 but
    are not relevant---lexically similar but semantically different.
    \item \textbf{Cross-encoder mined negatives:} Use a cross-encoder to find
    passages that a bi-encoder incorrectly scores highly.  Most effective but
    expensive.
\end{itemize}

\begin{keyinsight}[Sentence Embeddings Are the Backbone of RAG]
Modern Retrieval-Augmented Generation (RAG) systems depend entirely on the
quality of their embedding model.  The embedding model determines what gets
retrieved, and what gets retrieved determines the quality of the generated
answer.  When an interviewer asks about RAG, they are implicitly testing
whether you understand embedding quality, similarity search, and the
tradeoffs between different embedding approaches.
\end{keyinsight}


% ============================================================================
% SECTION 7: EMBEDDING EVALUATION
% ============================================================================
\section{Embedding Evaluation}
\label{sec:embedding_evaluation}

Evaluating embedding quality is non-trivial and is a common interview topic
because it tests whether a candidate thinks critically about measurement.

\subsection{Intrinsic Evaluation}

Intrinsic methods evaluate embeddings in isolation, without a downstream task.

\paragraph{Word Similarity.}
Compute cosine similarity between word pairs and correlate with human
judgments.  Standard datasets include:
\begin{itemize}
    \item \textbf{WordSim-353} (Finkelstein et al., 2001): 353 word pairs
    rated for similarity.
    \item \textbf{SimLex-999} (Hill et al., 2015): 999 pairs rated for
    \emph{genuine similarity} (not just relatedness---``car'' and ``road''
    are related but not similar).
    \item \textbf{MEN} (Bruni et al., 2014): 3,000 pairs rated for
    relatedness.
\end{itemize}
Evaluation metric: Spearman's rank correlation $\rho$ between model scores
and human judgments.

\paragraph{Analogy Tasks.}
Test whether the embedding space encodes relational structure.  The Google
analogy dataset contains 19,544 questions of the form ``$a$ is to $b$ as $c$
is to ?''  covering semantic (``Athens'' : ``Greece'' :: ``Oslo'' : ?) and
syntactic (``walk'' : ``walked'' :: ``run'' : ?) categories.

\paragraph{Clustering.}
Embed words and apply k-means or hierarchical clustering.  Inspect whether
resulting clusters correspond to meaningful categories (animals, countries,
verbs of motion).

\subsection{Extrinsic Evaluation}

Extrinsic methods evaluate embeddings by using them in a downstream task and
measuring task performance.  This is generally more informative than intrinsic
evaluation because it tests what actually matters---utility for a real
application.

\begin{warningbox}
Good intrinsic scores do not guarantee good extrinsic performance, and vice
versa.  A model that excels at word analogy tasks may perform poorly on NER
or sentiment analysis.  Always prefer extrinsic evaluation when the target
application is known.
\end{warningbox}

\subsection{MTEB: The Modern Benchmark}

The Massive Text Embedding Benchmark (Muennighoff et al., 2023) evaluates
sentence embeddings across 8 task categories and 58 datasets:
\begin{enumerate}
    \item Classification (e.g., sentiment, topic)
    \item Clustering (e.g., Reddit, arXiv)
    \item Pair Classification (e.g., duplicate detection)
    \item Re-ranking
    \item Retrieval (e.g., MS MARCO, NQ)
    \item Semantic Textual Similarity
    \item Summarization
    \item Bitext Mining (parallel sentence alignment)
\end{enumerate}

MTEB is now the standard benchmark for evaluating general-purpose embedding
models and is hosted as a public leaderboard on Hugging Face.

\subsection{Comparison of Embedding Methods}

\begin{table}[H]
\centering
\caption{Summary comparison of embedding methods across key dimensions.}
\label{tab:embedding_comparison}
\begin{tabularx}{\textwidth}{l c c c c c}
\toprule
\textbf{Method} & \textbf{Contextual} & \textbf{OOV} &
\textbf{GPU Needed} & \textbf{Typical $d$} & \textbf{Quality} \\
\midrule
One-hot       & No  & No  & No  & $V$ (sparse)  & Minimal \\
Word2Vec      & No  & No  & No  & 100--300      & Good \\
GloVe         & No  & No  & No  & 100--300      & Good \\
FastText      & No  & Yes & No  & 100--300      & Good+ \\
ELMo          & Yes & Yes & Yes & 512--1024     & Very Good \\
SBERT         & Yes & Yes & Yes & 384--768      & Excellent \\
E5/BGE/GTE    & Yes & Yes & Yes & 384--1024     & State-of-art \\
\bottomrule
\end{tabularx}
\end{table}


% ============================================================================
% SECTION 8: INTERVIEW QUESTIONS
% ============================================================================
\section{Interview Questions}
\label{sec:embedding_interview}

\begin{interviewq}
\textbf{Q: Explain the Word2Vec Skip-gram model. What is the training
objective? What is negative sampling and why is it needed?}

\badgeLfive{L5} \quad \badgetype{Mathematical} \quad \badgefreq{\freqhigh}

\stronganswer
Word2Vec Skip-gram learns word embeddings by training a model to predict
context words given a center word.  Each word $w$ has two vectors: an input
embedding $\vw_w$ and an output embedding $\vw'_w$.  Given center word $w_I$,
the probability of observing context word $w_O$ is:
\[
    P(w_O \given w_I)
    = \frac{\exp(\vw'_{w_O}{}^\top \vw_{w_I})}
           {\sum_{w=1}^{V} \exp(\vw'_w{}^\top \vw_{w_I})}
\]
The training objective maximizes the sum of $\log P(\text{context} \given
\text{center})$ over all (center, context) pairs in the corpus.

The problem is the denominator: it sums over the entire vocabulary $V$ (often
100K+ words) for every training example.  Negative sampling replaces this
expensive softmax with a binary classification task.  For each true
(center, context) pair, we sample $k$ noise words and train a binary
classifier to distinguish real pairs from noise.  The loss becomes:
\[
    \loss = -\log\sigma(\vw'_{w_O}{}^\top \vw_{w_I})
            - \sum_{i=1}^{k} \E_{w_i \sim P_n}\bigl[
              \log\sigma(-\vw'_{w_i}{}^\top \vw_{w_I})\bigr]
\]
The noise distribution is the unigram distribution raised to the 3/4 power,
which smooths the distribution and makes negative samples more informative.
This reduces training from $O(V)$ to $O(k)$ per example.

\redflags
\begin{itemize}
    \item Not knowing that Skip-gram has two sets of embeddings (input and
    output).
    \item Saying ``negative sampling samples wrong words'' without explaining
    the binary classification formulation.
    \item Being unable to explain why the 3/4 exponent on the noise
    distribution helps.
    \item Confusing Skip-gram (predict context from center) with CBOW
    (predict center from context).
\end{itemize}

\interviewtest{Mathematical depth, understanding of computational bottlenecks,
ability to explain optimization tricks and \emph{why} they work.}

\goodgreat
\begin{itemize}
    \item \badgeLfive{L5}: States the objective, explains negative sampling as
    binary classification, gives the noise distribution formula.
    \item \badgeLsix{L6}: Explains the connection to PMI (Levy and Goldberg,
    2014), discusses the effect of the number of negative samples, and
    mentions subsampling of frequent words.
    \item \badgeLseven{L7}: Derives that Skip-gram with negative sampling
    implicitly factorizes a shifted PMI matrix, connects to GloVe's
    explicit matrix factorization, discusses implications for embedding
    geometry and how this understanding informs modern contrastive learning
    objectives.
\end{itemize}

\followups{What is the connection between Word2Vec and PMI? How do you choose
the number of negative samples? What is subsampling and why does it help?
How does window size affect the type of relationships captured?}
\end{interviewq}

\vspace{0.5em}

% --- Question 2 ---
\begin{interviewq}
\textbf{Q: What are the differences between Word2Vec, GloVe, and FastText?
When would you choose each one?}

\badgeLfive{L5} \quad \badgetype{Conceptual} \quad \badgefreq{\freqhigh}

\stronganswer
All three produce static word embeddings but differ in their training approach
and capabilities.

\textbf{Word2Vec} (Mikolov et al., 2013) trains a shallow neural network on
local context windows.  Skip-gram predicts context from center word; CBOW
predicts center from context.  Training is online (streaming) and scales to
very large corpora.  Limitation: treats words as atomic units, no OOV
handling.

\textbf{GloVe} (Pennington et al., 2014) first computes a global
co-occurrence matrix, then optimizes a weighted least-squares objective so
that word vector dot products approximate the log co-occurrence count.  It
explicitly captures global statistics that Word2Vec's local windows might
miss, but requires pre-computing the co-occurrence matrix ($O(V^2)$ space,
though sparse in practice).

\textbf{FastText} (Bojanowski et al., 2017) extends Word2Vec by representing
each word as a bag of character n-grams.  A word's vector is the sum of its
n-gram vectors.  This enables OOV handling (decompose unseen words into known
n-grams), captures morphological structure, and produces better embeddings
for morphologically rich languages.

\textbf{When to choose each:}
\begin{itemize}
    \item \emph{Word2Vec:} General-purpose embeddings, very large corpora,
    when morphology is not a concern (e.g., English for well-curated text).
    \item \emph{GloVe:} When you want pre-trained vectors (Stanford provides
    high-quality pre-trained GloVe), or when you specifically need embeddings
    that capture global co-occurrence patterns.
    \item \emph{FastText:} Morphologically rich languages, noisy text
    (social media, user-generated content with typos), specialized domains
    with many OOV terms.  Also when you need fast supervised classification
    (FastText's supervised mode).
\end{itemize}

In practice, for most modern systems, the choice between these three matters
less than whether to use them at all versus a contextual model.  But for
low-resource, low-latency, or edge-deployment scenarios, FastText remains a
strong choice.

\redflags
\begin{itemize}
    \item Saying ``they are all the same, just use BERT.''
    \item Not knowing that GloVe uses a global co-occurrence matrix while
    Word2Vec uses local windows.
    \item Not knowing that FastText handles OOV via character n-grams.
    \item Being unable to articulate when static embeddings are preferred
    over contextual ones.
\end{itemize}

\interviewtest{Breadth of knowledge across embedding methods, understanding
of tradeoffs, ability to make principled engineering decisions.}

\goodgreat
\begin{itemize}
    \item \badgeLfive{L5}: Describes all three methods correctly, gives a
    reasonable decision framework.
    \item \badgeLsix{L6}: Discusses the theoretical connection (all implicitly
    factorize co-occurrence statistics), mentions Levy and Goldberg's
    unification result, and discusses production considerations (latency,
    model size, OOV rates).
    \item \badgeLseven{L7}: Compares with modern subword tokenizers (BPE
    handles some of the same OOV issues FastText addresses), discusses
    hybrid approaches (FastText features in conjunction with contextual
    models), and explains how static embeddings are still used as components
    of larger systems.
\end{itemize}

\followups{How does Levy and Goldberg's result connect Word2Vec to SVD on
a PMI matrix? When would you use static embeddings in a modern system?
How does FastText's approach differ from BPE tokenization?}
\end{interviewq}

\vspace{0.5em}

% --- Question 3 ---
\begin{interviewq}
\textbf{Q: Why are contextual embeddings better than static embeddings?
Give a concrete example of when static embeddings fail.}

\badgeLfive{L5} \quad \badgetype{Conceptual} \quad \badgefreq{\freqhigh}

\stronganswer
Static embeddings (Word2Vec, GloVe, FastText) assign a single fixed vector to
each word regardless of context.  Contextual embeddings (ELMo, BERT) assign
\emph{different} vectors depending on the surrounding words.

The fundamental failure mode of static embeddings is \textbf{polysemy}.
Consider the word ``crane'':
\begin{itemize}
    \item ``The \textbf{crane} lifted the steel beam.'' (construction
    equipment)
    \item ``We observed a \textbf{crane} at the lake.'' (bird)
    \item ``She had to \textbf{crane} her neck to see.'' (verb: to stretch)
\end{itemize}
A static embedding places all three senses at a single point---a compromise
that accurately represents none of them.  A contextual model produces three
distinct vectors, each appropriate for its context.

Beyond polysemy, contextual embeddings capture:
\begin{itemize}
    \item \textbf{Word sense disambiguation} (automatically, without explicit
    sense inventories).
    \item \textbf{Syntactic role} (``running'' as noun vs.\ verb gets
    different representations).
    \item \textbf{Sentence-level semantics} (negation, quantification,
    conditionals affect word representations).
\end{itemize}

Quantitatively, the shift from static to contextual embeddings produced
massive improvements: BERT improved GLUE benchmark scores by 7+ points over
ELMo (which itself improved significantly over static embeddings).

\redflags
\begin{itemize}
    \item Giving only ``contextual is better because BERT is better'' without
    explaining the polysemy problem.
    \item Not being able to give a concrete example.
    \item Saying static embeddings are ``useless'' (they are still useful
    for efficiency, baselines, and as components of larger systems).
\end{itemize}

\interviewtest{Understanding of \emph{why} the paradigm shifted, not just
\emph{that} it did.  Ability to reason about failure modes and design
implications.}

\goodgreat
\begin{itemize}
    \item \badgeLfive{L5}: Explains polysemy with a clear example,
    mentions ELMo and BERT as contextual models.
    \item \badgeLsix{L6}: Discusses what different layers capture (syntax
    vs.\ semantics in ELMo/BERT), explains the practical tradeoff (static
    embeddings are fast, contextual requires GPU inference), and mentions
    scenarios where static embeddings still win.
    \item \badgeLseven{L7}: Discusses the emergence of word senses in
    BERT's embedding space (Wiedemann et al., 2019), explains how
    contextual embeddings enable transfer learning (pre-train once, adapt to
    many tasks), and connects to the broader pre-training paradigm.
\end{itemize}

\followups{What does ELMo capture at different layers? How does BERT create
contextual embeddings differently from ELMo? When would you still use static
embeddings in a modern system?}
\end{interviewq}

\vspace{0.5em}

% --- Question 4 ---
\begin{interviewq}
\textbf{Q: How would you evaluate the quality of word or sentence embeddings?}

\badgeLfive{L5} \quad \badgetype{Conceptual} \quad \badgefreq{\freqmed}

\stronganswer
Embedding evaluation has two categories: \textbf{intrinsic} and
\textbf{extrinsic}.

\textbf{Intrinsic evaluation} measures embedding properties in isolation:
\begin{itemize}
    \item \emph{Word similarity:} Compute cosine similarity for word pairs
    and measure Spearman correlation with human ratings (WordSim-353,
    SimLex-999).
    \item \emph{Analogy tasks:} Test $\vw_a - \vw_b + \vw_c \approx \vw_d$
    (Google analogy dataset).
    \item \emph{Clustering:} Check whether embeddings cluster by meaningful
    categories.
\end{itemize}

\textbf{Extrinsic evaluation} uses embeddings in a downstream task and
measures task metrics---this is generally more informative because it tests
what actually matters.  For sentence embeddings, the MTEB benchmark
evaluates across 8 task categories and 58 datasets.

\textbf{Critically}, good intrinsic scores do not guarantee good extrinsic
performance.  An embedding space with perfect analogy structure might still
perform poorly on NER or retrieval.  I would always prioritize evaluation
on the target task---if I am building a retrieval system, I evaluate with
retrieval metrics (Recall@K, NDCG), not word similarity scores.

\redflags
\begin{itemize}
    \item Only mentioning intrinsic evaluation.
    \item Not knowing specific benchmarks (WordSim-353, SimLex-999, MTEB).
    \item Saying ``just look at the nearest neighbors'' without quantitative
    metrics.
\end{itemize}

\interviewtest{Critical thinking about evaluation, awareness that proxies can
mislead, practical measurement discipline.}

\goodgreat
\begin{itemize}
    \item \badgeLfive{L5}: Describes both intrinsic and extrinsic evaluation
    with specific examples.
    \item \badgeLsix{L6}: Discusses the disconnect between intrinsic and
    extrinsic metrics, mentions MTEB, and explains how to set up a
    task-specific evaluation pipeline.
    \item \badgeLseven{L7}: Discusses embedding space geometry (anisotropy,
    uniformity---Ethayarajh, 2019), explains how to detect and correct
    representation degeneration, and designs an evaluation strategy that
    combines multiple signals for high-stakes deployment decisions.
\end{itemize}

\followups{What is the MTEB benchmark? What is embedding anisotropy and why
is it a problem? How would you compare two embedding models for a production
search system?}
\end{interviewq}

\vspace{0.5em}

% --- Question 5 ---
\begin{interviewq}
\textbf{Q: What is the distributional hypothesis and why is it foundational
to all word embedding methods?}

\badgeLfive{L5} \quad \badgetype{Conceptual} \quad \badgefreq{\freqmed}

\stronganswer
The distributional hypothesis, articulated by J.~R.~Firth in 1957, states
that ``a word is characterized by the company it keeps''---words appearing in
similar contexts tend to have similar meanings.

This is the theoretical foundation for \emph{every} word embedding method:
\begin{itemize}
    \item \emph{Count-based methods} (PMI, SVD) directly model co-occurrence
    statistics in context windows.
    \item \emph{Word2Vec} trains to predict context words, implicitly
    learning to encode distributional patterns.
    \item \emph{GloVe} optimizes over the global co-occurrence matrix.
    \item \emph{Contextual models} (ELMo, BERT) predict masked or next tokens
    from context, learning representations that capture the full range of
    distributional patterns.
\end{itemize}

The hypothesis operationalizes ``meaning'' as something computable from
observable data (text corpora) rather than requiring formal semantic
definitions.  This is why unsupervised training on raw text works: the
distributional signal IS the semantic signal.

The hypothesis has limitations: distributional similarity does not capture all
aspects of meaning.  ``Dog'' and ``cat'' are distributionally similar
(similar contexts) but semantically distinct.  Antonyms like ``hot'' and
``cold'' appear in nearly identical contexts but have opposite meanings.
Grounded or experiential meaning (what ``red'' looks like) cannot be
captured from text distribution alone.

\redflags
\begin{itemize}
    \item Not attributing the hypothesis or knowing where it comes from.
    \item Treating it as a curiosity rather than the foundation of the field.
    \item Not acknowledging its limitations (antonyms, grounding).
\end{itemize}

\interviewtest{Depth of understanding of \emph{why} embeddings work, not just
\emph{how}.  Awareness of theoretical foundations.}

\goodgreat
\begin{itemize}
    \item \badgeLfive{L5}: States the hypothesis correctly, connects it to
    at least two embedding methods.
    \item \badgeLsix{L6}: Discusses limitations (antonyms, grounding),
    explains how all embedding methods operationalize the hypothesis
    differently.
    \item \badgeLseven{L7}: Discusses Harris (1954) and Firth (1957) in
    historical context, connects to modern probing studies that show
    contextual models capture distributional patterns beyond simple
    co-occurrence, and discusses the ``meaning-from-text'' debate
    (Bender and Koller, 2020).
\end{itemize}

\followups{Can distributional methods capture antonym relationships? What
aspects of meaning are NOT captured by distributional methods? How does the
distributional hypothesis relate to BERT's masked language modeling?}
\end{interviewq}

\vspace{0.5em}

% --- Question 6 ---
\begin{interviewq}
\textbf{Q: How does FastText handle out-of-vocabulary words? Why is this
an improvement over Word2Vec?}

\badgeLfive{L5} \quad \badgetype{Conceptual} \quad \badgefreq{\freqmed}

\stronganswer
FastText represents each word as the sum of its character n-gram embeddings
(typically $n = 3$ to $6$) plus a whole-word embedding.  For example,
``where'' with $n=3$ is decomposed into the n-grams
$\{\texttt{<wh}, \texttt{whe}, \texttt{her}, \texttt{ere},
\texttt{re>}\}$ plus the whole-word token $\texttt{<where>}$.

For an OOV word, FastText decomposes it into character n-grams and sums the
vectors of those n-grams that were observed during training.  Since the
n-gram vocabulary is far smaller than the word vocabulary and covers
almost all possible character sequences, FastText virtually eliminates the
OOV problem.

This is a significant improvement over Word2Vec, which has zero representation
for OOV words---the model simply cannot process them.  This improvement
matters most for:
\begin{itemize}
    \item Morphologically rich languages where compound words and
    agglutination create enormous vocabularies.
    \item Noisy text with typos and misspellings.
    \item Domain-specific terminology with productive morphology.
\end{itemize}

Additionally, FastText produces better embeddings even for in-vocabulary
words, because morphologically related words (``happy,'' ``unhappy,''
``happiness'') share n-grams and therefore have related vectors.

\redflags
\begin{itemize}
    \item Saying ``FastText just uses character-level embeddings'' (it uses
    n-grams, not individual characters).
    \item Not knowing that the word vector is the \emph{sum} of n-gram
    vectors.
    \item Confusing FastText's subword approach with BPE tokenization (they
    solve different problems in different ways).
\end{itemize}

\interviewtest{Understanding of practical engineering considerations in NLP,
awareness of the OOV problem and its solutions.}

\goodgreat
\begin{itemize}
    \item \badgeLfive{L5}: Explains the n-gram decomposition and summation
    correctly, gives an example.
    \item \badgeLsix{L6}: Discusses trade-offs (larger n-gram vocabulary,
    increased memory), compares with BPE tokenization (FastText is at the
    embedding level; BPE is at the tokenization level), and gives
    language-specific examples.
    \item \badgeLseven{L7}: Discusses how modern subword tokenizers (BPE,
    SentencePiece) have partly superseded FastText's approach to OOV
    handling, and explains the remaining niche where FastText's approach
    is advantageous (e.g., extremely fast inference without neural network
    forward pass, cross-lingual alignment via shared n-grams).
\end{itemize}

\followups{How does FastText's approach compare to BPE tokenization for
handling OOV? When would you choose FastText over a transformer-based model
in production? How does the n-gram range affect quality?}
\end{interviewq}

\vspace{0.5em}

% --- Question 7 ---
\begin{interviewq}
\textbf{Q: Explain SimCSE. How does using dropout as data augmentation
create meaningful positive pairs for contrastive learning?}

\badgeLsix{L6} \quad \badgetype{Conceptual} \quad \badgefreq{\freqmed}

\stronganswer
SimCSE (Gao et al., 2021) is a contrastive learning framework for sentence
embeddings.  The unsupervised version uses a remarkably simple trick: pass the
same sentence through the encoder \emph{twice}, with different dropout masks
applied in each forward pass.  The two resulting embeddings form a
\emph{positive pair}.  Other sentences in the batch serve as negatives.  The
model is trained with an InfoNCE loss to bring positive pairs closer and push
negatives apart.

\textbf{Why dropout works as augmentation.}  Standard dropout randomly zeroes
a fraction of hidden units during each forward pass.  When the same input
passes through the network twice, dropout produces two \emph{slightly
different} internal representations---the same semantic content processed
through two different ``sub-networks'' of the transformer.  The key insight is
that this is a \emph{minimal perturbation} that preserves semantics: unlike
text-level augmentations (synonym replacement, back-translation, word
deletion), dropout cannot change the meaning of the input because the
\emph{input is identical} in both passes.

The supervised variant uses NLI data: entailment pairs as positives and
contradiction pairs as hard negatives, consistently outperforming the
unsupervised version.

SimCSE also addressed a key problem in sentence embedding spaces:
\emph{anisotropy}---the tendency for BERT sentence embeddings to occupy a
narrow cone in the vector space, making cosine similarity unreliable.
Contrastive learning produces more uniform, isotropic representations.

\redflags
\begin{itemize}
    \item Not understanding that the positive pair uses the \emph{same}
    input sentence (just different dropout masks).
    \item Conflating SimCSE with generic data augmentation approaches.
    \item Not knowing about the anisotropy problem that SimCSE addresses.
\end{itemize}

\interviewtest{Understanding of contrastive learning, creative approaches to
data augmentation, and practical knowledge of sentence embedding methods.}

\goodgreat
\begin{itemize}
    \item \badgeLfive{L5}: Explains the mechanism correctly, states that
    dropout creates the positive pairs.
    \item \badgeLsix{L6}: Explains why dropout works (minimal semantic
    perturbation, different ``sub-networks''), discusses the
    anisotropy problem, and compares unsupervised vs.\ supervised
    variants.
    \item \badgeLseven{L7}: Discusses the connection to the
    uniformity-alignment framework (Wang and Isola, 2020), explains how
    contrastive learning corrects representation degeneration, and connects
    to modern embedding training pipelines that build on SimCSE's
    insights (E5, BGE use similar contrastive objectives at scale).
\end{itemize}

\followups{What is the anisotropy problem in BERT embeddings? How does
contrastive learning produce more isotropic representations? How does
SimCSE compare to SBERT? What is the uniformity-alignment framework?}
\end{interviewq}

\vspace{0.5em}

% --- Question 8 ---
\begin{interviewq}
\textbf{Q: How would you choose an embedding model for a production semantic
search system?}

\badgeLsix{L6} \quad \badgetype{System Design} \quad \badgefreq{\freqmed}

\stronganswer
Choosing an embedding model for production semantic search requires balancing
quality, latency, cost, and operational complexity.  I would evaluate along
these dimensions:

\textbf{1.\ Quality on the target domain.}  Evaluate candidate models on
a held-out set of queries and relevant documents from the target domain.
MTEB scores give a general signal, but domain-specific evaluation is
essential.  An embedding model trained on academic text may perform poorly
on e-commerce queries.

\textbf{2.\ Latency budget.}  Embedding models range from FastText
(sub-millisecond CPU inference) to large transformer encoders (10--50ms
on GPU).  For real-time search with strict latency SLAs, smaller models
(MiniLM, E5-small) or pre-computed embeddings with ANN search are necessary.

\textbf{3.\ Dimensionality and storage.}  Each document embedding must be
stored in the vector index.  For 100M documents at 768 dimensions in
float32, that is $\sim$300GB.  Matryoshka embeddings allow truncation to
128 or 256 dimensions with modest quality loss, reducing storage by 3--6x.
Quantization (scalar, product quantization) further reduces storage.

\textbf{4.\ Multilingual requirements.}  If the system serves multiple
languages, a multilingual model (E5-multilingual, BGE-m3) is needed.  Be
aware of the ``tokenizer tax''---non-Latin scripts produce more tokens,
reducing effective context length and increasing latency.

\textbf{5.\ Fine-tuning potential.}  If you have domain-specific labeled
data (query-document relevance pairs), fine-tuning a general embedding model
typically improves quality by 5--15\%.  Consider whether the model supports
efficient fine-tuning (LoRA, continued contrastive training).

\textbf{6.\ Hybrid architecture.}  In most production systems, dense
embeddings are combined with sparse retrieval (BM25) in a hybrid
architecture.  The embedding model handles semantic matching; BM25 handles
exact lexical matching.  A cross-encoder re-ranker then scores the top
candidates.  The embedding model should be evaluated as part of this full
pipeline, not in isolation.

\redflags
\begin{itemize}
    \item Choosing based solely on MTEB leaderboard ranking without
    domain-specific evaluation.
    \item Ignoring latency, storage, and cost constraints.
    \item Not mentioning hybrid retrieval (dense + sparse).
    \item Proposing a cross-encoder for first-stage retrieval (does not
    scale).
\end{itemize}

\interviewtest{Systems thinking, production awareness, ability to balance
competing constraints, understanding of the full retrieval pipeline.}

\goodgreat
\begin{itemize}
    \item \badgeLfive{L5}: Identifies key dimensions (quality, latency, cost)
    and names specific models.
    \item \badgeLsix{L6}: Discusses hybrid retrieval, Matryoshka embeddings,
    domain-specific fine-tuning, and the full retrieval pipeline including
    re-ranking.
    \item \badgeLseven{L7}: Discusses quantization strategies for the vector
    index, multi-stage retrieval design (coarse ANN $\to$ re-rank $\to$
    cross-encoder), monitoring embedding quality drift over time, and the
    decision framework for when to fine-tune vs.\ use off-the-shelf models
    including cost-benefit analysis.
\end{itemize}

\followups{How would you handle embedding drift as your data distribution
changes? How would you evaluate retrieval quality end-to-end? What is the
tradeoff between Matryoshka dimension truncation and quality?}
\end{interviewq}

\vspace{0.5em}

% --- Question 9 ---
\begin{interviewq}
\textbf{Q: Explain the connection between Word2Vec, GloVe, and matrix
factorization on co-occurrence statistics. Why does this connection matter?}

\badgeLsix{L6} \quad \badgetype{Mathematical} \quad \badgefreq{\freqmed}

\stronganswer
Levy and Goldberg (2014) proved a remarkable result: Word2Vec Skip-gram with
negative sampling is implicitly factorizing a shifted pointwise mutual
information (PMI) matrix.  Specifically, the inner product of the learned
embeddings converges to:
\[
    \vw_i^\top \vw'_j = \text{PMI}(i, j) - \log k
\]
where $k$ is the number of negative samples.  This means Skip-gram is
performing a weighted, low-rank factorization of the word-context PMI matrix.

GloVe, by contrast, \emph{explicitly} optimizes a weighted factorization of
the log co-occurrence matrix:
$\vw_i^\top \vw'_j + b_i + b'_j \approx \log X_{ij}$.

The classical SVD approach directly factorizes the PPMI matrix using truncated
SVD.

All three methods are fundamentally doing the same thing: factorizing a
matrix derived from co-occurrence statistics into low-dimensional word
vectors.  They differ in the weighting scheme, the specific matrix being
factorized (PPMI vs.\ log counts vs.\ implicitly shifted PMI), and the
optimization procedure (SGD on local windows vs.\ batch optimization on the
full matrix vs.\ SVD).

This connection matters because it: (1)~explains why all three methods produce
similar-quality embeddings despite very different formulations; (2)~shows that
``predictive'' and ``count-based'' are not fundamentally different paradigms;
and (3)~provides theoretical grounding for why these methods work---they
capture the statistical structure of language as encoded in co-occurrence
patterns.

\redflags
\begin{itemize}
    \item Treating Word2Vec as ``completely different'' from count-based
    methods.
    \item Not knowing the Levy and Goldberg result.
    \item Being unable to state what matrix Word2Vec implicitly factorizes.
\end{itemize}

\interviewtest{Depth of theoretical understanding, ability to see unifying
principles across seemingly different methods, research literacy.}

\goodgreat
\begin{itemize}
    \item \badgeLfive{L5}: States that Word2Vec and GloVe both relate to
    co-occurrence statistics.
    \item \badgeLsix{L6}: Cites Levy and Goldberg, states the specific
    shifted PMI result, and explains how GloVe's objective relates.
    \item \badgeLseven{L7}: Discusses the implications for hyperparameter
    choices (window size controls the co-occurrence matrix, negative
    samples control the PMI shift), connects to modern contrastive learning
    (which also performs implicit spectral decomposition), and discusses
    how this understanding informed the design of subsequent methods.
\end{itemize}

\followups{What matrix does Word2Vec implicitly factorize? How does the
number of negative samples affect this factorization? Does this unification
extend to contextual embeddings like BERT?}
\end{interviewq}

\vspace{0.5em}

% --- Question 10 ---
\begin{interviewq}
\textbf{Q: Walk me through how you would build a text similarity system
for a product with 50 million documents, starting from embedding selection
to serving.}

\badgeLsix{L6} \quad \badgetype{System Design} \quad \badgefreq{\freqmed}

\stronganswer
I would design this as a multi-stage pipeline.

\textbf{Stage 1: Embedding Model Selection.}  Evaluate candidate models
(E5-large, BGE-large, a domain-fine-tuned variant) on a representative sample
of document pairs with similarity labels.  Measure Spearman correlation for
similarity and Recall@K for retrieval.  Choose based on quality-latency
tradeoff---for 50M documents, a model producing 384--768 dimensional
embeddings is practical.  If quality is critical, use Matryoshka-trained
embeddings to enable dimension truncation for candidate generation.

\textbf{Stage 2: Offline Embedding Generation.}  Embed all 50M documents in a
batch pipeline.  At 768 dimensions in float32, the index is approximately
150GB.  Apply scalar quantization (float32 $\to$ int8) to reduce to
$\sim$37GB.  Store in a vector database (Pinecone, Weaviate, Qdrant, or
self-hosted FAISS/ScaNN).

\textbf{Stage 3: Index Construction.}  Build an ANN (Approximate Nearest
Neighbor) index.  HNSW is the standard choice for production systems: good
recall at sub-millisecond latency.  IVF-PQ is an alternative when memory is
constrained.  Tune recall-latency tradeoff on the evaluation set.

\textbf{Stage 4: Query Pipeline.}  At query time: embed the query (10--30ms
on GPU), search the ANN index for top-100 candidates ($<$5ms), then re-rank
with a cross-encoder for the top-20 ($\sim$50ms).  Total latency: $<$100ms.

\textbf{Stage 5: Monitoring.}  Track embedding quality over time.  Set up
golden evaluation sets, monitor retrieval metrics weekly, and detect
distribution drift.  Periodically refresh embeddings as the document corpus
evolves.

\redflags
\begin{itemize}
    \item Proposing brute-force cosine similarity over 50M documents
    (does not scale).
    \item Not mentioning ANN search or vector databases.
    \item Ignoring the re-ranking stage.
    \item Not considering storage and quantization.
\end{itemize}

\interviewtest{End-to-end systems thinking, production awareness, ability to
design a complete pipeline with concrete numbers for latency, storage, and
throughput.}

\goodgreat
\begin{itemize}
    \item \badgeLfive{L5}: Describes embedding + ANN search + re-ranking at
    a high level.
    \item \badgeLsix{L6}: Provides concrete latency and storage estimates,
    discusses quantization, names specific ANN algorithms and vector
    databases, and includes a hybrid retrieval component.
    \item \badgeLseven{L7}: Designs for incremental index updates, discusses
    sharding and replication for availability, addresses cold-start for new
    documents, considers A/B testing frameworks for comparing embedding
    models, and designs the monitoring and retraining pipeline.
\end{itemize}

\followups{How would you handle incremental updates when new documents
arrive? How would you A/B test two different embedding models in production?
What ANN algorithm would you choose and why?}
\end{interviewq}
